\documentclass[final,a4paper,12pt]{article}
\usepackage{notes}

\title{Geophysics III}
\author{Steady-State Geotherms}
\date{}

\begin{document}
\maketitle
%\input{../psintro}

\noindent This problem set is worth a total of 20 points to the total problem set score.

\begin{enumerate}
    %\clearpage
    \item \textbf{[4 pts.] Thermal conductivity} \\
Using the mineralogical composition chart and the thermal conductivities
listed in the table below, compute the thermal conductivity for

\par\medskip
\begin{figure}[hbt]
  \centering{
    \includegraphics[angle=90]{figures/igneous_rock_minerals.eps}
  }
  \caption{Mineralogical compositions of igneous rocks.}
  \label{fig:geo}
\end{figure}
\medskip

\par\medskip
\begin{table}[hbt]
    \centering
    \begin{threeparttable}
    \caption{Thermal conductivity of common materials.}
    \label{tab:tc}
    \begin{tabular}{@{}lc@{}}
        \toprule
        {} &
        $k$ \\
        mineral &
        \si{W.m^{-1}.K^{-1}} \\
        \midrule
        $\alpha$-quartz           & 8.8 \\
        orthoclase                & 1.8 \\
        plagioclase, felsic       & 1.8 \\
        plagioclase, mafic        & 1.6 \\
        mica, muscovite           & 2.1 \\
        mica, biotite             & 1.9 \\
        hornblende (amphibole)    & 2.7 \\
        clinopyroxene             & 4.3 \\
        orthopyroxene             & 3.4 \\
        olivine                   & 4.8 \\
        spinel                    & 9 \\
        garnet                    & 3.9 \\
        \bottomrule
    \end{tabular}
    \begin{tablenotes}
        \item All conductivities estimated at \SI{25}{^\circ C}.
    \end{tablenotes}
    \end{threeparttable}
\end{table}
\medskip

\begin{enumerate}
    \item granite
    \item diorite
    \item gabbro, assuming all pyroxene is clinopyroxene.
    \item peridotite

    Phase changes associated with pressure turn plagioclase into spinel and
    then into garnet.  Therefore, compute the thermal conductivity of
    peridotite with a composition of 67\% olivine, 17\%
    orthopyroxene, 8\% clinopyroxene, and 8\% garnet.
\end{enumerate}
\FloatBarrier



    %clearpage
    \item {\bf [10 pts.] Continental Geotherms} \\
    This problem is completed with the python-based codes, \verb+ss_isostasy.py+, which you can find on the course GitHub page. Run in VSCode or python.

    {\it On paper:} 
    \begin{enumerate}
        \item Let's start by computing a geotherm (temperature as a function of depth) roughly consistent with a craton.  We will keep it relatively simple by dividing the lithosphere into three layers: an upper crust; lower crust; and mantle lithosphere (Figure~\ref{fig:layers}).

        \par\medskip
        \begin{figure}[hbtp]
            \centering
            \includegraphics[]{figures/layers.eps}
            \caption{A three layered geotherm model.}
            \label{fig:layers}
        \end{figure}
        \medskip

        Using the geotherm solution in Appendix~\ref{apx:g}, compute the temperature at each of the interfaces above (except the lithosphere--asthenosphere boundary).  Use $q_0 = \SI{40}{mW.m^{-2}}$ and $A_{1} = \SI{1.0}{\mu W.m^{-3}}$.  For the lithosphere--asthenosphere boundary, assume the convecting mantle temperature is 1300\degC.  This geotherm will be our reference for all future elevation computations.

        Plot your result and interpolate between layers.  {\bf What should the shape of your curve be?}
    \end{enumerate}

    {\it On computer:}\\
    Complete the rest using \verb+ss_isostasy.py+, a python calculator that will compute geotherms,
    \begin{enumerate}
        \item  Now, computing the reference geotherm roughly consistent with a craton, using the same layering as above (Figure~\ref{fig:layers}).  Entering the appropriate properties into the table and text boxes.
        
        \item What causes the kinks in the curve, especially at the Moho?
        
        \item Now let's examine the sensitivity of geotherms to variations in heat production.  This test will illustrate why reducing uncertainties in crustal heat production are crucial for properly estimating lithospheric temperatures.

        Compute three geotherms: (1) increase upper crustal heat production by 25\%; (2) increase lower crustal heat production by 25\%; and (3) increase mantle lithospheric heat production by 25\%.  All other properties remain the same as those in Figure~\ref{fig:layers}.

        {\bf Does increasing heat production increase or decrease the lithospheric temperatures?  How much does each test affect the temperature at the Moho?  How about the depth at which the adiabat is intersected?}

        \item Now let's compute a family of geotherms with heat flow values of 40, 60, 80, 100~mW~m$^{-2}$.  The upper crustal heat production and surface heat flow are the only value we will vary.  To determine the heat production we will assume a partition coefficient of $P = 0.75$.  The upper crustal heat production can be computed by
        \begin{equation}
            A_{1} = (1 - P) \frac{q_0}{D},
        \end{equation}
        where $D$ is the thickness of the upper crust, 10~km.

        Using your combinations of $q_0$ and $A_1$ above, compute each geotherm.  Aside from $q_0$ and $A_1$, all other inputs to the model should be the same.

        \item Compute the elevations of each geotherm from the geotherm family above (the equations it uses are found in Appendix~\ref{apx:tis}).  Use the 40~mW~m$^{-2}$ geotherm as the reference.  Assume the thermal expansivity, $\alpha$, is $3 \times 10^{-5}$~\si{K^{-1}}.  Hint: you won't need to compute the elevation for the 40~mW~m$^{-2}$ geotherm since the reference elevation is defined as 0~km.  Plot the elevations as a function of surface heat flow.  To compute the total elevation, you will need to compute the individual elevation contributions for four steps, one for each lithospheric layer and one for depths below the intersection with the mantle adiabat.

        {\bf How much is the elevation affected by heat flow over the range of typical continental heat flows (40 to \SI{100}{mW.m^{-2}})?}

        \item Now lets see if we can estimate the heat production of the Australian provinces.  Start by computing another set of geotherms with heat flow $q_0 = \SI{80}{mW.m^{-2}}$ with $P = (0.75, 0.55, 0.45, 0.35)$.

        Repeat the procedure from 3 for each new P.  Plot the new geotherm family.

        \item Compute the elevation differences for geotherms with $q_0 = \SI{80}{mW.m^{-2}}$.  Use the $q_0 = 40$~mW~m$^{2}$ geotherm as the reference.

        Repeat the steps used to produce the geotherm family and elevations above with the new heat production parameters.  Plot each result.

        {\bf Which value of heat production is most consistent with the reference cratonic geotherm?  Could this be the explanation for our elevation anomalies?}

    \end{enumerate}


    %clearpage
    \item \label{itm:rad} {\bf [6 pts.] Thermal stability of cratons} \\
Cratons are typically viewed as regions that are thermally cool, which
contributes to their longevity/stability.  In this problem, you will explore the
effect of crustal radioactivity on crustal temperatures both for the present-day
Earth and at a time \SI{3}{Ga} when many of these stable regions were being first
created.

This problem is accompanied by a spreadsheet (see template online) that you will
need to complete as part of the exercises below.
\begin{enumerate}
    \item Compute the average heat production for a granitic and gabbroic
        compositions at the present day for crystallization ages of $<$\SI{200}{Ma},
        from 2000 to \SI{2800}{Ma}, and prior to \SI{2800}{Ma}.  The formulae should be
        directly entered into the worksheet labeled "heat production" and copied
        and modified as necessary for the other two rock ages.
        
        Note: the concentration of potassium is provided as K$_2$O, so you will
        need to convert to K.  And both K and U have multiple isotopes, so you
        will need to account for that in your calculation using the constants
        provided in the worksheet.

    \item In the same sheet, compute the average heat production for these same
        rock types and ages adjusted to a time 3 billion years ago (\SI{3}{Ga}).

    \item The values you previously computed will be automatically copied into
        the worksheet labeled ``temperature''.  We want to compute a set of
        geotherms in this new sheet for a simplified crust created for six
        crustal heat production scenarios corresponding to the different age
        rocks at both present day and at \SI{3}{Ga}.

        We will make some simplifying assumptions for a typically moderate
        crustal heat flow:
        \begin{itemize}
            \item upper crust is \SI{16}{km} thick; 
            \item lower crust is \SI{24}{km} thick;
            \item basal heat flow is \SI{36}{mW.m^{-2}}---use the conservation of
                energy to compute the heat flow at the top of each layer, all
                the way back to the surface; and
            \item thermal conductivity is temperature variable---the equations
                for thermal conductivity are already computed for you and will
                automatically update once you have computed the temperature.
        \end{itemize}

        If you enter the equations correctly into the first empty heat flow and
        temperature cells, you should be able to copy them to the end of each
        column.  With the appropriate \$ on cell columns/rows, you will also be
        able to copy the heat flow and temperature columns from the first
        scenario to all others.

    \item The remainder of the problem leads up to computing heat production for
        three hypothetical regions at a point \SI{3}{Ga} in the past so that the
        geotherms in the past.

        Start by determining the relationship between half-life and radioactive
        decay constants.  The concentration of a radioelement, P(t), is given by
        \begin{equation}
            P(t) = P_0 \exp (-\lambda t),
        \end{equation}
        where $P_0$ is the initial concentration, and $\lambda$ is the decay
        constant.  Hint: a half-life is defined as the time at which the ratio
        of the radioelement concentration is one-half the initial concentration.
        Show your work.

    \item Compute the heat production of for the heat production concentrations
        at a point in time \SI{3}{Ga} in the past.  To compute the heat production in
        the past, you will need to adjust the concentrations of the
        radioisotopes.

    \item Compute one-dimensional, steady-state geotherms for each of the
        situations above to the Moho at a time \SI{3}{Ga} in the past.

    \item Explain, in light of your computations, how changes in heat production
        concentrations of igneous rocks helps lead to stability of Archean
        lithosphere.
\end{enumerate}
\FloatBarrier

\end{enumerate}

\appendix
\section{Computing geotherms}\label{apx:g}

The solution for the temperature at the base of the first layer,
$T_{1}$, can then be estimated using the homogeneous solution,
\begin{equation}
  T_1 = T_0 + \frac{q_0}{k_1} (z_1 - z_0) -
  \frac{A_1}{2 k_1} (z_1 - z_0)^2.
\end{equation}
The heat flow at the base of the first layer is equal to
\begin{equation}
    q_{1} = q_0 - A_1 (z_1 - z_0).
\end{equation}
The homogeneous solution can then be used to estimate the temperature at
the base of the second layer using, $T_1$, as the temperature at the
top and $q_1$ as the heat flow.  The temperature of all subsequent
layers can be found using the same method boot-strapping down through
the domain of interest.  

Solving for the temperature at the base of the $i$-th layer,
\begin{equation}
  T_{i+1} = T_i + \frac{q_i}{k_{i+1}} (z_{i+1} - z_i)
  - \frac{A_{i+1}}{2 k_{i+1}} (z_{i+1} - z_i)^2.
\end{equation}
It is also necessary to know the heat loss at the base of each layer in
order to compute temperature.  Since heat flow is determined by the heat
flowing into a layer plus the heat generated within, the heat flow at
the bottom of each layer is given by
\begin{equation}
    q_{i+1} = q_i - A_{i+1} (z_{i+1} - z_i).
\end{equation}
This solution forms the basis of most 1-D steady-state geotherms.


\section{Thermal isostasy}\label{apx:tis}
It can be shown that the relative elevation difference between a region with
geotherm $T(z)$ and a reference region $T_{\text{ref}}$ is given by
\begin{equation}
    \varepsilon = \alpha \int_{0}^{z_{\text{max}}} [T(z) - T_{\text{ref}}] \,dz.
\end{equation}
For our piecewise geotherms, we can compute an elevation within each depth
interval by
\begin{equation}
    \varepsilon_i = \alpha \int_{z_i}^{z_{i+1}} [T(z) - T_{\text{ref}}] \,dz,
\end{equation}
and summing the individual steps
\begin{equation}
    \varepsilon = \sum_{i = 1}^N \varepsilon_i.
\end{equation}
Within each step, we can integrate difference in temperature by integrating each
geotherm separately and taking the difference since the integral operator is
linear,
\begin{equation}
    \varepsilon_i = \alpha \left[
    \int_{z_i}^{z_{i+1}} T(z) \,dz - \int_{z_i}^{z_{i+1}} T_{\text{ref}} \,dz
    \right].
\end{equation}
The integral of a single geotherm is given by
\begin{align}
    \int_{z_i}^{z_{i+1}} T(z) \,dz & = \int_{z_i}^{z_{i+1}}
        \left[ T_{i} 
        + \frac{q_i}{k_{i+1}} (z - z_i)
        - \frac{A_{i+1}}{2 k_{i+1}} (z - z_i)^2
        \right] \,dz, \\
    & = T_{i} z \left.
        + \frac{q_i}{2 k_{i+1}} (z - z_i)^2
        - \frac{A_{i+1}}{6 k_{i+1}} (z - z_i)^3
        \right|_{z_i}^{z_{i+1}}, \\
    & = T_{i} (z_{i+1} - z_i) 
        + \frac{q_i}{2 k_{i+1}} (z_{i+1} - z_i)^2
        - \frac{A_{i+1}}{6 k_{i+1}} (z_{i+1} - z_i)^3 .
\end{align}

\end{document}